\section{样本协方差在高维为何会失真}
\label{sec:11-Mathematical-Statistics/09-Random-Matrix-Theory/01}

\subsection{当你信任的特征值突然散了架}

做 PCA 时你得到一个协方差矩阵，算出特征值，画出 scree plot。如果在低维数据（比如 10 个变量、1000 个样本）上，这张图通常很稳定：前几个特征值突出，后面整齐地衰减，变点明显。现在你面对一组基因表达数据——2000 个基因、200 个样本。真实协方差矩阵没人知道，但你发现样本特征值的分布从极小的正数一路拖到天文数字，和直觉中的``大部分方差集中在前几个方向''差得很远。

是数据质量出了问题，还是代码写错了？都不是。

做一个干净的模拟来建立直觉。设 \(X_1,\ldots,X_n\in\R^p\) 是独立同分布样本，真实协方差就是单位阵 \(\Sigma=I\)——每个坐标独立、方差为 1，绝对零结构，主成分分析在这里没有有意义的方向，每个方向方差完全相同。样本协方差

\[
\widehat\Sigma=\frac1n\sum_{i=1}^n(X_i-\bar X)(X_i-\bar X)^\trans
\]

在低维渐近（\(p\) 固定、\(n\to\infty\)）下收敛到 \(\Sigma\)，每个特征值都一致地逼近真值 1。这个结论是数理统计教科书的基石——大数定律保证样本协方差逐元素逼近真实协方差。经典的结论是 \(\widehat\Sigma\to\Sigma\) 几乎必然。

但如果 \(p\) 与 \(n\) 同阶增长——记 \(p/n\to c>0\)——这个结论整体失效。取极端情况 \(p=n\)（\(c=1\)），样本协方差的最小特征值可以接近 0，最大特征值可以接近 4。谱从集中在 \(\{1\}\) 这个单点，弥散成一个连续的区间 \([0,4]\)。你自己跑一遍就信了：用 Python 生成 \(n=p=500\) 的标准正态数据，算样本协方差的特征值，画直方图——一个光滑的曲线从 0 延伸到 4，峰值在 1 附近，但两侧拖出了不可忽视的质量。增大 \(n\) 到 5000 但保持 \(p/n=1\)，直方图形状几乎不变——它不是有限样本的小波动，而是高维渐近的真实极限。你的样本量增加了十倍，特征值的弥散却毫发未缩。

\subsection{无偏估计为什么保不住特征值}

这里有一个容易忽略的直觉断层：\(\widehat\Sigma\) 确实是 \(\Sigma\) 的无偏估计——逐元素取期望，每个 \(\widehat\Sigma_{ij}\) 的期望等于 \(\Sigma_{ij}\)。在低维 regime 里，逐元素无偏就足够了——大数定律会收拾一切。但特征值是矩阵元素的非线性函数，无偏性不穿越非线性变换。一百个各自无偏的元素的随机涨落，经过特征分解这个高度非线性的工序后，不会碰巧互相抵消，而是系统性地推动特征值向外弥散——就像一百个单独无偏的测量，在算出它们的平方和之后，得到的不是无偏的方差估计（需要除以 \(n-1\) 而不是 \(n\)，这就是 Bessel 校正的由来）。非线性变换把无偏性吃掉了。

更几何地看：\(\widehat\Sigma\) 的秩至多为 \(n\)。你用一个至多 \(n\) 维的子空间去承载 \(p\) 个方向上的方差信息。若 \(p>n\)，至少 \(p-n\) 个特征值精确为零——样本协方差奇异，逆不存在，任何需要 \(\widehat\Sigma^{-1}\) 的方法直接崩掉。即使 \(p=n\)，秩刚好够用，最小的几个特征方向也几乎没有独立的数据支撑——只有个别样本恰好落在了那个方向上，随机涨落把这些方向的特征值压向零；最大的几个特征方向吸收了本该属于其他方向的方差，膨胀到远超 1。

秩的限制还带来一个直接后果：特征值的平均仍然是 1（因为 trace 的无偏性得以保留），所以大特征值的偏大和小特征值的偏小刚好在平均值意义上互相补偿。谱的整体中心是对的，但弥散幅度完全失控。

这与``病态问题''有本质区别。病态指真实 \(\Sigma\) 本身条件数巨大，问题是内在的——即使知道真实的 \(\Sigma\)，求逆依然数值不稳定。这里的现象是：即使真实 \(\Sigma\) 完美良态（单位阵，条件数为 1），估计量的谱也会因为高维而系统扭曲。样本量增大不能解决它——只要 \(p/n\) 保持非零常数，失真就不消失。

从参数计数看：样本协方差要估计 \(p(p+1)/2\) 个自由参数（对称矩阵上三角），却只有 \(np\) 个标量观测。当 \(p\sim n\)，每个参数分到的有效观测数趋于常数——约 \(2n/(p+1)\approx 2/c\)。每个参数只分到少量观测，逐元素涨落没有机会被大数定律抹平。

\subsection{两种渐近世界观：低维 vs 高维}

低维渐近（\(p\) 固定、\(n\to\infty\)）给出的结论是：
\begin{itemize}
\tightlist
\item \(\widehat\Sigma\to\Sigma\) 几乎必然；
\item 每个特征值一致收敛到真值；
\item 标准误按 \(1/\sqrt n\) 缩小；
\item \(n\) 足够大时，可以忽略估计误差；
\item 高维只是把低维公式里填上更大的 \(p\)。
\end{itemize}

高维渐近（\(p/n\to c\in(0,\infty)\)）的结论则是：
\begin{itemize}
\tightlist
\item \(\widehat\Sigma\) 的特征值不收敛到 \(\Sigma\) 的特征值；
\item 而是收敛到一个由 \(c\) 和 \(\Sigma\) 的谱共同决定的确定分布；
\item 即使 \(n\to\infty\)，样本特征值也不会坍缩回真值；
\item 经典标准误公式在 \(c>0\) 时是不一致的——\(n\) 越大它越错。
\end{itemize}

\(c\) 是控制失真程度的关键参数。\(c=0.01\) 时失真很轻，经典公式近似可用——每个参数约分到 200 个观测。\(c=0.1\) 时最大特征值已可感知地偏大。\(c=0.5\) 时每个参数只分到约 4 个观测，失真明显。\(c=1\) 时谱宽拉到 \([0,4]\)，每个参数约 2 个观测。\(c=2\)（\(p=2n\)）时样本协方差奇异，至少一半特征值为零——无论真协方差多良态。

统计推断必须声明自己在哪个 regime 下工作。用低维标准误公式在高维数据上算置信区间，等价于假设每个参数有几百个观测而实际上不到一个——数值看起来没问题（程序不报错），但覆盖概率完全错误。Part 11 前面的章节默认低维或中等维度的世界观；本单元处理的是 \(p\) 与 \(n\) 同阶的 regime。

\subsection{失真的下游代价：从 PCA 到投资组合}

样本协方差的谱失真不是一个孤立的技术细节。它会逐层渗透到所有吃协方差估计的方法：

\begin{itemize}
\tightlist
\item
  \textbf{PCA}：前几个样本特征值系统性偏大——最大特征方向吸收了本该属于其他方向的方差。你可能错误地判断前三个主成分解释了 80\% 方差，而实际只有 30\%。样本特征向量与真实主方向之间可能出现不可忽略的夹角，即使真实特征值之间有明显差距。用有偏的特征值和有偏的特征向量做降维，再在降维空间做下游推断——失真被带入后续每一步。
\item
  \textbf{马氏距离与判别分析}：\(\widehat\Sigma^{-1}\) 把最小特征方向的微小值倒数化放大成巨大权重。在真实分布下平平无奇的样本，在估计的马氏距离下可能被判为极端离群。LDA 和 QDA 的分类边界被拖向最小特征方向，高维中的分类稳定度远不如低维。
\item
  \textbf{协方差正则化}：Ledoit--Wolf 收缩、阈值化、低秩加稀疏分解等方法的设计前提，正是承认样本协方差的谱不可靠、需要向某个先验结构（单位阵、对角阵或因子模型）收缩。不知失真来源与形状，无据可依地选择收缩目标，修正可能比不修更差。
\item
  \textbf{Markowitz 投资组合}：最优权重公式含 \(\widehat\Sigma^{-1}\)。谱失真让组合权重对输入极度敏感——昨天和今天的协方差估计只有微小差异，最优组合却可能从全仓 A 变成全仓 B。这就是``估计误差主导优化误差''的来源：不是你的优化算错了，是你的协方差估计在高维下不再可信。
\end{itemize}

下一节给出失真的精确形状：Marchenko--Pastur 律。当你看到那个分布从随机矩阵的特征值直方图中浮现出来时，低维直觉和高维现实之间的落差会变成可精确计算的对象——而不是一团模糊的``高维不可信''的担忧。
